% =====================================================================
% TikZ macros for drawing pretzel knot diagrams.
%
% USAGE: % =====================================================================
% TikZ macros for drawing pretzel knot diagrams.
%
% USAGE: % =====================================================================
% TikZ macros for drawing pretzel knot diagrams.
%
% USAGE: % =====================================================================
% TikZ macros for drawing pretzel knot diagrams.
%
% USAGE: \input{tikz_pretzel_macros} in your dissertation preamble.
%
% Requires the following packages and libraries to be loaded:
%   \usepackage{tikz}
%   \usepackage{xcolor}
%   \usetikzlibrary{calc}
%
% Defines:
%   \crossing{x}{y}{w}{h}{sign}{color}    -- single crossing
%   \tangle{xc}{ytop}{ybot}{k}{color}     -- vertical tangle of k crossings
%   \pretzel{p1}{p2}{p3}{color}           -- 3-tangle pretzel P(p1,p2,p3)
%   \bandhighlight                        -- dashed pink box on third tangle
%
% These also define the colours \thesisblue, \thesisred, \bandfill
% (used in the dissertation's banner/highlight scheme).
% =====================================================================

% --- Colours ---------------------------------------------------------
\providecolor{thesisblue}{HTML}{1F4E79}
\providecolor{thesisred}{HTML}{C00000}
\providecolor{bandfill}{HTML}{FFE5E5}

% ---------------------------------------------------------------------
% \crossing{x}{y}{w}{h}{sign}{color}
%   A single crossing centred at (x, y), of width w and height h.
%   sign > 0: right-handed (SW->NE strand over NW->SE strand).
%   sign < 0: left-handed  (NW->SE strand over SW->NE strand).
% ---------------------------------------------------------------------
\newcommand{\crossing}[6]{%
  \pgfmathsetmacro{\eps}{0.08*#3}%
  \pgfmathsetmacro{\hw}{0.5*#3}%
  \pgfmathsetmacro{\hh}{0.5*#4}%
  \ifnum#5>0\relax
    \draw[#6, line width=1.1pt, line cap=round]
      ({#1-\hw},{#2+\hh}) -- ({#1-\eps},{#2+\eps*#4/#3});
    \draw[#6, line width=1.1pt, line cap=round]
      ({#1+\eps},{#2-\eps*#4/#3}) -- ({#1+\hw},{#2-\hh});
    \draw[#6, line width=1.1pt, line cap=round]
      ({#1-\hw},{#2-\hh}) -- ({#1+\hw},{#2+\hh});
  \else
    \draw[#6, line width=1.1pt, line cap=round]
      ({#1-\hw},{#2-\hh}) -- ({#1-\eps},{#2-\eps*#4/#3});
    \draw[#6, line width=1.1pt, line cap=round]
      ({#1+\eps},{#2+\eps*#4/#3}) -- ({#1+\hw},{#2+\hh});
    \draw[#6, line width=1.1pt, line cap=round]
      ({#1-\hw},{#2+\hh}) -- ({#1+\hw},{#2-\hh});
  \fi
}

% ---------------------------------------------------------------------
% \tangle{xc}{ytop}{ybot}{k}{color}
%   A vertical tangle of |k| crossings of sign(k) between y = ytop
%   and y = ybot, centred horizontally at x = xc. k = 0 gives two
%   parallel vertical strands.
% ---------------------------------------------------------------------
\newcommand{\tangle}[5]{%
  \pgfmathsetmacro{\tanglew}{0.55}%
  \pgfmathsetmacro{\hw}{0.5*\tanglew}%
  \ifnum#4=0\relax
    \draw[#5, line width=1.1pt, line cap=round]
      ({#1-\hw},#2) -- ({#1-\hw},#3);
    \draw[#5, line width=1.1pt, line cap=round]
      ({#1+\hw},#2) -- ({#1+\hw},#3);
  \else
    \ifnum#4>0\relax
      \def\sgnval{1}%
      \pgfmathtruncatemacro{\absk}{#4}%
    \else
      \def\sgnval{-1}%
      \pgfmathtruncatemacro{\absk}{-(#4)}%
    \fi
    \pgfmathsetmacro{\totalh}{#2-(#3)}%
    \pgfmathsetmacro{\hper}{\totalh/\absk}%
    \foreach \i in {1,...,\absk} {%
      \pgfmathsetmacro{\cy}{#2-(\i-0.5)*\hper}%
      \crossing{#1}{\cy}{\tanglew}{\hper*0.85}{\sgnval}{#5}%
    }%
  \fi
}

% ---------------------------------------------------------------------
% \pretzel{p1}{p2}{p3}{color}
%   The standard diagram of the 3-tangle pretzel P(p1, p2, p3),
%   centred at (0, 0) and approximately 7 x 8 units in size.
% ---------------------------------------------------------------------
\newcommand{\pretzel}[4]{%
  \pgfmathsetmacro{\tanglew}{0.55}%
  \pgfmathsetmacro{\hw}{0.5*\tanglew}%
  \def\ytop{1.6}%
  \def\ybot{-1.6}%
  \def\xleft{-2.0}%
  \def\xmid{0.0}%
  \def\xright{2.0}%
  \tangle{\xleft}{\ytop}{\ybot}{#1}{#4}%
  \tangle{\xmid}{\ytop}{\ybot}{#2}{#4}%
  \tangle{\xright}{\ytop}{\ybot}{#3}{#4}%
  \draw[#4, line width=1.1pt, line cap=round]
    ({\xleft+\hw}, \ytop)
      arc[start angle=180, end angle=0, x radius=0.725, y radius=0.32];
  \draw[#4, line width=1.1pt, line cap=round]
    ({\xmid+\hw}, \ytop)
      arc[start angle=180, end angle=0, x radius=0.725, y radius=0.32];
  \draw[#4, line width=1.1pt, line cap=round]
    ({\xleft+\hw}, \ybot)
      arc[start angle=-180, end angle=0, x radius=0.725, y radius=0.32];
  \draw[#4, line width=1.1pt, line cap=round]
    ({\xmid+\hw}, \ybot)
      arc[start angle=-180, end angle=0, x radius=0.725, y radius=0.32];
  \draw[#4, line width=1.1pt, line cap=round]
    ({\xleft-\hw}, \ytop) -- ({\xleft-\hw}, {\ytop+0.4})
      arc[start angle=180, end angle=0,
          x radius={(\xright-\xleft)/2+\hw}, y radius=0.85]
      -- ({\xright+\hw}, \ytop);
  \draw[#4, line width=1.1pt, line cap=round]
    ({\xleft-\hw}, \ybot) -- ({\xleft-\hw}, {\ybot-0.4})
      arc[start angle=-180, end angle=0,
          x radius={(\xright-\xleft)/2+\hw}, y radius=0.85]
      -- ({\xright+\hw}, \ybot);
  \node[anchor=north, font=\bfseries, fill=white, inner sep=2pt]
    at (\xleft,  {\ybot-1.4}) {$#1$};
  \node[anchor=north, font=\bfseries, fill=white, inner sep=2pt]
    at (\xmid,   {\ybot-1.4}) {$#2$};
  \node[anchor=north, font=\bfseries, fill=white, inner sep=2pt]
    at (\xright, {\ybot-1.4}) {$#3$};
}

% ---------------------------------------------------------------------
% \bandhighlight
%   A dashed pink box around the third tangle of \pretzel, highlighting
%   it as the parametric "twisting band" of the Hedden-Watson construction.
% ---------------------------------------------------------------------
\newcommand{\bandhighlight}{%
  \draw[thesisred, dashed, line width=0.8pt,
        fill=bandfill, fill opacity=0.35]
    (1.55, -1.75) rectangle (2.45, 1.75);
}

% ---------------------------------------------------------------------
% \trefoilcs{s1}{s2}{color}
%   The standard diagram of the connect-sum T(s1) # T(s2) of two
%   trefoils, where s1, s2 in {+1, -1} indicate the chirality of each
%   summand (+1 = right-handed, -1 = left-handed).
%
%   Drawn as two side-by-side closed 2-braids joined by the trivial
%   connect-sum band. Centred at (0, 0).
% ---------------------------------------------------------------------
\newcommand{\trefoilcs}[3]{%
  \pgfmathsetmacro{\tanglew}{0.55}%
  \pgfmathsetmacro{\hw}{0.5*\tanglew}%
  \def\ytop{1.5}%
  \def\ybot{-1.5}%
  \def\xL{-1.5}%
  \def\xR{1.5}%
  % Pre-evaluate the crossing counts (3 with the sign of s1, s2)
  \pgfmathtruncatemacro{\trefL}{(#1)*3}%
  \pgfmathtruncatemacro{\trefR}{(#2)*3}%
  % Two trefoil tangles (each = 3 crossings of given chirality)
  \tangle{\xL}{\ytop}{\ybot}{\trefL}{#3}%
  \tangle{\xR}{\ytop}{\ybot}{\trefR}{#3}%
  % Closure of left trefoil: top arc from left-of-tangle going up-and-over to
  % right-of-tangle. This is the "outside" closure of the 2-braid.
  \draw[#3, line width=1.1pt, line cap=round]
    ({\xL-\hw}, \ytop) -- ({\xL-\hw}, {\ytop+0.35})
      arc[start angle=180, end angle=0, x radius=\hw, y radius=0.32]
      -- ({\xL+\hw}, \ytop);
  % Closure of right trefoil: same on the right side
  \draw[#3, line width=1.1pt, line cap=round]
    ({\xR-\hw}, \ytop) -- ({\xR-\hw}, {\ytop+0.35})
      arc[start angle=180, end angle=0, x radius=\hw, y radius=0.32]
      -- ({\xR+\hw}, \ytop);
  % Connect-sum band at the bottom: trivial band joining the two trefoils.
  % It loops from the bottom-right of the left tangle to the bottom-left
  % of the right tangle, dipping down. Then we close off the two outer
  % bottom strands with a large arc going around.
  \draw[#3, line width=1.1pt, line cap=round]
    ({\xL+\hw}, \ybot)
      arc[start angle=180, end angle=360, x radius={(\xR-\xL-\tanglew)/2}, y radius=0.55];
  % Outer bottom closure: leftmost-bottom to rightmost-bottom around the underside.
  \draw[#3, line width=1.1pt, line cap=round]
    ({\xL-\hw}, \ybot) -- ({\xL-\hw}, {\ybot-0.45})
      arc[start angle=180, end angle=360, x radius={(\xR-\xL)/2+\hw}, y radius=0.85]
      -- ({\xR+\hw}, \ybot);
}
 in your dissertation preamble.
%
% Requires the following packages and libraries to be loaded:
%   \usepackage{tikz}
%   \usepackage{xcolor}
%   \usetikzlibrary{calc}
%
% Defines:
%   \crossing{x}{y}{w}{h}{sign}{color}    -- single crossing
%   \tangle{xc}{ytop}{ybot}{k}{color}     -- vertical tangle of k crossings
%   \pretzel{p1}{p2}{p3}{color}           -- 3-tangle pretzel P(p1,p2,p3)
%   \bandhighlight                        -- dashed pink box on third tangle
%
% These also define the colours \thesisblue, \thesisred, \bandfill
% (used in the dissertation's banner/highlight scheme).
% =====================================================================

% --- Colours ---------------------------------------------------------
\providecolor{thesisblue}{HTML}{1F4E79}
\providecolor{thesisred}{HTML}{C00000}
\providecolor{bandfill}{HTML}{FFE5E5}

% ---------------------------------------------------------------------
% \crossing{x}{y}{w}{h}{sign}{color}
%   A single crossing centred at (x, y), of width w and height h.
%   sign > 0: right-handed (SW->NE strand over NW->SE strand).
%   sign < 0: left-handed  (NW->SE strand over SW->NE strand).
% ---------------------------------------------------------------------
\newcommand{\crossing}[6]{%
  \pgfmathsetmacro{\eps}{0.08*#3}%
  \pgfmathsetmacro{\hw}{0.5*#3}%
  \pgfmathsetmacro{\hh}{0.5*#4}%
  \ifnum#5>0\relax
    \draw[#6, line width=1.1pt, line cap=round]
      ({#1-\hw},{#2+\hh}) -- ({#1-\eps},{#2+\eps*#4/#3});
    \draw[#6, line width=1.1pt, line cap=round]
      ({#1+\eps},{#2-\eps*#4/#3}) -- ({#1+\hw},{#2-\hh});
    \draw[#6, line width=1.1pt, line cap=round]
      ({#1-\hw},{#2-\hh}) -- ({#1+\hw},{#2+\hh});
  \else
    \draw[#6, line width=1.1pt, line cap=round]
      ({#1-\hw},{#2-\hh}) -- ({#1-\eps},{#2-\eps*#4/#3});
    \draw[#6, line width=1.1pt, line cap=round]
      ({#1+\eps},{#2+\eps*#4/#3}) -- ({#1+\hw},{#2+\hh});
    \draw[#6, line width=1.1pt, line cap=round]
      ({#1-\hw},{#2+\hh}) -- ({#1+\hw},{#2-\hh});
  \fi
}

% ---------------------------------------------------------------------
% \tangle{xc}{ytop}{ybot}{k}{color}
%   A vertical tangle of |k| crossings of sign(k) between y = ytop
%   and y = ybot, centred horizontally at x = xc. k = 0 gives two
%   parallel vertical strands.
% ---------------------------------------------------------------------
\newcommand{\tangle}[5]{%
  \pgfmathsetmacro{\tanglew}{0.55}%
  \pgfmathsetmacro{\hw}{0.5*\tanglew}%
  \ifnum#4=0\relax
    \draw[#5, line width=1.1pt, line cap=round]
      ({#1-\hw},#2) -- ({#1-\hw},#3);
    \draw[#5, line width=1.1pt, line cap=round]
      ({#1+\hw},#2) -- ({#1+\hw},#3);
  \else
    \ifnum#4>0\relax
      \def\sgnval{1}%
      \pgfmathtruncatemacro{\absk}{#4}%
    \else
      \def\sgnval{-1}%
      \pgfmathtruncatemacro{\absk}{-(#4)}%
    \fi
    \pgfmathsetmacro{\totalh}{#2-(#3)}%
    \pgfmathsetmacro{\hper}{\totalh/\absk}%
    \foreach \i in {1,...,\absk} {%
      \pgfmathsetmacro{\cy}{#2-(\i-0.5)*\hper}%
      \crossing{#1}{\cy}{\tanglew}{\hper*0.85}{\sgnval}{#5}%
    }%
  \fi
}

% ---------------------------------------------------------------------
% \pretzel{p1}{p2}{p3}{color}
%   The standard diagram of the 3-tangle pretzel P(p1, p2, p3),
%   centred at (0, 0) and approximately 7 x 8 units in size.
% ---------------------------------------------------------------------
\newcommand{\pretzel}[4]{%
  \pgfmathsetmacro{\tanglew}{0.55}%
  \pgfmathsetmacro{\hw}{0.5*\tanglew}%
  \def\ytop{1.6}%
  \def\ybot{-1.6}%
  \def\xleft{-2.0}%
  \def\xmid{0.0}%
  \def\xright{2.0}%
  \tangle{\xleft}{\ytop}{\ybot}{#1}{#4}%
  \tangle{\xmid}{\ytop}{\ybot}{#2}{#4}%
  \tangle{\xright}{\ytop}{\ybot}{#3}{#4}%
  \draw[#4, line width=1.1pt, line cap=round]
    ({\xleft+\hw}, \ytop)
      arc[start angle=180, end angle=0, x radius=0.725, y radius=0.32];
  \draw[#4, line width=1.1pt, line cap=round]
    ({\xmid+\hw}, \ytop)
      arc[start angle=180, end angle=0, x radius=0.725, y radius=0.32];
  \draw[#4, line width=1.1pt, line cap=round]
    ({\xleft+\hw}, \ybot)
      arc[start angle=-180, end angle=0, x radius=0.725, y radius=0.32];
  \draw[#4, line width=1.1pt, line cap=round]
    ({\xmid+\hw}, \ybot)
      arc[start angle=-180, end angle=0, x radius=0.725, y radius=0.32];
  \draw[#4, line width=1.1pt, line cap=round]
    ({\xleft-\hw}, \ytop) -- ({\xleft-\hw}, {\ytop+0.4})
      arc[start angle=180, end angle=0,
          x radius={(\xright-\xleft)/2+\hw}, y radius=0.85]
      -- ({\xright+\hw}, \ytop);
  \draw[#4, line width=1.1pt, line cap=round]
    ({\xleft-\hw}, \ybot) -- ({\xleft-\hw}, {\ybot-0.4})
      arc[start angle=-180, end angle=0,
          x radius={(\xright-\xleft)/2+\hw}, y radius=0.85]
      -- ({\xright+\hw}, \ybot);
  \node[anchor=north, font=\bfseries, fill=white, inner sep=2pt]
    at (\xleft,  {\ybot-1.4}) {$#1$};
  \node[anchor=north, font=\bfseries, fill=white, inner sep=2pt]
    at (\xmid,   {\ybot-1.4}) {$#2$};
  \node[anchor=north, font=\bfseries, fill=white, inner sep=2pt]
    at (\xright, {\ybot-1.4}) {$#3$};
}

% ---------------------------------------------------------------------
% \bandhighlight
%   A dashed pink box around the third tangle of \pretzel, highlighting
%   it as the parametric "twisting band" of the Hedden-Watson construction.
% ---------------------------------------------------------------------
\newcommand{\bandhighlight}{%
  \draw[thesisred, dashed, line width=0.8pt,
        fill=bandfill, fill opacity=0.35]
    (1.55, -1.75) rectangle (2.45, 1.75);
}

% ---------------------------------------------------------------------
% \trefoilcs{s1}{s2}{color}
%   The standard diagram of the connect-sum T(s1) # T(s2) of two
%   trefoils, where s1, s2 in {+1, -1} indicate the chirality of each
%   summand (+1 = right-handed, -1 = left-handed).
%
%   Drawn as two side-by-side closed 2-braids joined by the trivial
%   connect-sum band. Centred at (0, 0).
% ---------------------------------------------------------------------
\newcommand{\trefoilcs}[3]{%
  \pgfmathsetmacro{\tanglew}{0.55}%
  \pgfmathsetmacro{\hw}{0.5*\tanglew}%
  \def\ytop{1.5}%
  \def\ybot{-1.5}%
  \def\xL{-1.5}%
  \def\xR{1.5}%
  % Pre-evaluate the crossing counts (3 with the sign of s1, s2)
  \pgfmathtruncatemacro{\trefL}{(#1)*3}%
  \pgfmathtruncatemacro{\trefR}{(#2)*3}%
  % Two trefoil tangles (each = 3 crossings of given chirality)
  \tangle{\xL}{\ytop}{\ybot}{\trefL}{#3}%
  \tangle{\xR}{\ytop}{\ybot}{\trefR}{#3}%
  % Closure of left trefoil: top arc from left-of-tangle going up-and-over to
  % right-of-tangle. This is the "outside" closure of the 2-braid.
  \draw[#3, line width=1.1pt, line cap=round]
    ({\xL-\hw}, \ytop) -- ({\xL-\hw}, {\ytop+0.35})
      arc[start angle=180, end angle=0, x radius=\hw, y radius=0.32]
      -- ({\xL+\hw}, \ytop);
  % Closure of right trefoil: same on the right side
  \draw[#3, line width=1.1pt, line cap=round]
    ({\xR-\hw}, \ytop) -- ({\xR-\hw}, {\ytop+0.35})
      arc[start angle=180, end angle=0, x radius=\hw, y radius=0.32]
      -- ({\xR+\hw}, \ytop);
  % Connect-sum band at the bottom: trivial band joining the two trefoils.
  % It loops from the bottom-right of the left tangle to the bottom-left
  % of the right tangle, dipping down. Then we close off the two outer
  % bottom strands with a large arc going around.
  \draw[#3, line width=1.1pt, line cap=round]
    ({\xL+\hw}, \ybot)
      arc[start angle=180, end angle=360, x radius={(\xR-\xL-\tanglew)/2}, y radius=0.55];
  % Outer bottom closure: leftmost-bottom to rightmost-bottom around the underside.
  \draw[#3, line width=1.1pt, line cap=round]
    ({\xL-\hw}, \ybot) -- ({\xL-\hw}, {\ybot-0.45})
      arc[start angle=180, end angle=360, x radius={(\xR-\xL)/2+\hw}, y radius=0.85]
      -- ({\xR+\hw}, \ybot);
}
 in your dissertation preamble.
%
% Requires the following packages and libraries to be loaded:
%   \usepackage{tikz}
%   \usepackage{xcolor}
%   \usetikzlibrary{calc}
%
% Defines:
%   \crossing{x}{y}{w}{h}{sign}{color}    -- single crossing
%   \tangle{xc}{ytop}{ybot}{k}{color}     -- vertical tangle of k crossings
%   \pretzel{p1}{p2}{p3}{color}           -- 3-tangle pretzel P(p1,p2,p3)
%   \bandhighlight                        -- dashed pink box on third tangle
%
% These also define the colours \thesisblue, \thesisred, \bandfill
% (used in the dissertation's banner/highlight scheme).
% =====================================================================

% --- Colours ---------------------------------------------------------
\providecolor{thesisblue}{HTML}{1F4E79}
\providecolor{thesisred}{HTML}{C00000}
\providecolor{bandfill}{HTML}{FFE5E5}

% ---------------------------------------------------------------------
% \crossing{x}{y}{w}{h}{sign}{color}
%   A single crossing centred at (x, y), of width w and height h.
%   sign > 0: right-handed (SW->NE strand over NW->SE strand).
%   sign < 0: left-handed  (NW->SE strand over SW->NE strand).
% ---------------------------------------------------------------------
\newcommand{\crossing}[6]{%
  \pgfmathsetmacro{\eps}{0.08*#3}%
  \pgfmathsetmacro{\hw}{0.5*#3}%
  \pgfmathsetmacro{\hh}{0.5*#4}%
  \ifnum#5>0\relax
    \draw[#6, line width=1.1pt, line cap=round]
      ({#1-\hw},{#2+\hh}) -- ({#1-\eps},{#2+\eps*#4/#3});
    \draw[#6, line width=1.1pt, line cap=round]
      ({#1+\eps},{#2-\eps*#4/#3}) -- ({#1+\hw},{#2-\hh});
    \draw[#6, line width=1.1pt, line cap=round]
      ({#1-\hw},{#2-\hh}) -- ({#1+\hw},{#2+\hh});
  \else
    \draw[#6, line width=1.1pt, line cap=round]
      ({#1-\hw},{#2-\hh}) -- ({#1-\eps},{#2-\eps*#4/#3});
    \draw[#6, line width=1.1pt, line cap=round]
      ({#1+\eps},{#2+\eps*#4/#3}) -- ({#1+\hw},{#2+\hh});
    \draw[#6, line width=1.1pt, line cap=round]
      ({#1-\hw},{#2+\hh}) -- ({#1+\hw},{#2-\hh});
  \fi
}

% ---------------------------------------------------------------------
% \tangle{xc}{ytop}{ybot}{k}{color}
%   A vertical tangle of |k| crossings of sign(k) between y = ytop
%   and y = ybot, centred horizontally at x = xc. k = 0 gives two
%   parallel vertical strands.
% ---------------------------------------------------------------------
\newcommand{\tangle}[5]{%
  \pgfmathsetmacro{\tanglew}{0.55}%
  \pgfmathsetmacro{\hw}{0.5*\tanglew}%
  \ifnum#4=0\relax
    \draw[#5, line width=1.1pt, line cap=round]
      ({#1-\hw},#2) -- ({#1-\hw},#3);
    \draw[#5, line width=1.1pt, line cap=round]
      ({#1+\hw},#2) -- ({#1+\hw},#3);
  \else
    \ifnum#4>0\relax
      \def\sgnval{1}%
      \pgfmathtruncatemacro{\absk}{#4}%
    \else
      \def\sgnval{-1}%
      \pgfmathtruncatemacro{\absk}{-(#4)}%
    \fi
    \pgfmathsetmacro{\totalh}{#2-(#3)}%
    \pgfmathsetmacro{\hper}{\totalh/\absk}%
    \foreach \i in {1,...,\absk} {%
      \pgfmathsetmacro{\cy}{#2-(\i-0.5)*\hper}%
      \crossing{#1}{\cy}{\tanglew}{\hper*0.85}{\sgnval}{#5}%
    }%
  \fi
}

% ---------------------------------------------------------------------
% \pretzel{p1}{p2}{p3}{color}
%   The standard diagram of the 3-tangle pretzel P(p1, p2, p3),
%   centred at (0, 0) and approximately 7 x 8 units in size.
% ---------------------------------------------------------------------
\newcommand{\pretzel}[4]{%
  \pgfmathsetmacro{\tanglew}{0.55}%
  \pgfmathsetmacro{\hw}{0.5*\tanglew}%
  \def\ytop{1.6}%
  \def\ybot{-1.6}%
  \def\xleft{-2.0}%
  \def\xmid{0.0}%
  \def\xright{2.0}%
  \tangle{\xleft}{\ytop}{\ybot}{#1}{#4}%
  \tangle{\xmid}{\ytop}{\ybot}{#2}{#4}%
  \tangle{\xright}{\ytop}{\ybot}{#3}{#4}%
  \draw[#4, line width=1.1pt, line cap=round]
    ({\xleft+\hw}, \ytop)
      arc[start angle=180, end angle=0, x radius=0.725, y radius=0.32];
  \draw[#4, line width=1.1pt, line cap=round]
    ({\xmid+\hw}, \ytop)
      arc[start angle=180, end angle=0, x radius=0.725, y radius=0.32];
  \draw[#4, line width=1.1pt, line cap=round]
    ({\xleft+\hw}, \ybot)
      arc[start angle=-180, end angle=0, x radius=0.725, y radius=0.32];
  \draw[#4, line width=1.1pt, line cap=round]
    ({\xmid+\hw}, \ybot)
      arc[start angle=-180, end angle=0, x radius=0.725, y radius=0.32];
  \draw[#4, line width=1.1pt, line cap=round]
    ({\xleft-\hw}, \ytop) -- ({\xleft-\hw}, {\ytop+0.4})
      arc[start angle=180, end angle=0,
          x radius={(\xright-\xleft)/2+\hw}, y radius=0.85]
      -- ({\xright+\hw}, \ytop);
  \draw[#4, line width=1.1pt, line cap=round]
    ({\xleft-\hw}, \ybot) -- ({\xleft-\hw}, {\ybot-0.4})
      arc[start angle=-180, end angle=0,
          x radius={(\xright-\xleft)/2+\hw}, y radius=0.85]
      -- ({\xright+\hw}, \ybot);
  \node[anchor=north, font=\bfseries, fill=white, inner sep=2pt]
    at (\xleft,  {\ybot-1.4}) {$#1$};
  \node[anchor=north, font=\bfseries, fill=white, inner sep=2pt]
    at (\xmid,   {\ybot-1.4}) {$#2$};
  \node[anchor=north, font=\bfseries, fill=white, inner sep=2pt]
    at (\xright, {\ybot-1.4}) {$#3$};
}

% ---------------------------------------------------------------------
% \bandhighlight
%   A dashed pink box around the third tangle of \pretzel, highlighting
%   it as the parametric "twisting band" of the Hedden-Watson construction.
% ---------------------------------------------------------------------
\newcommand{\bandhighlight}{%
  \draw[thesisred, dashed, line width=0.8pt,
        fill=bandfill, fill opacity=0.35]
    (1.55, -1.75) rectangle (2.45, 1.75);
}

% ---------------------------------------------------------------------
% \trefoilcs{s1}{s2}{color}
%   The standard diagram of the connect-sum T(s1) # T(s2) of two
%   trefoils, where s1, s2 in {+1, -1} indicate the chirality of each
%   summand (+1 = right-handed, -1 = left-handed).
%
%   Drawn as two side-by-side closed 2-braids joined by the trivial
%   connect-sum band. Centred at (0, 0).
% ---------------------------------------------------------------------
\newcommand{\trefoilcs}[3]{%
  \pgfmathsetmacro{\tanglew}{0.55}%
  \pgfmathsetmacro{\hw}{0.5*\tanglew}%
  \def\ytop{1.5}%
  \def\ybot{-1.5}%
  \def\xL{-1.5}%
  \def\xR{1.5}%
  % Pre-evaluate the crossing counts (3 with the sign of s1, s2)
  \pgfmathtruncatemacro{\trefL}{(#1)*3}%
  \pgfmathtruncatemacro{\trefR}{(#2)*3}%
  % Two trefoil tangles (each = 3 crossings of given chirality)
  \tangle{\xL}{\ytop}{\ybot}{\trefL}{#3}%
  \tangle{\xR}{\ytop}{\ybot}{\trefR}{#3}%
  % Closure of left trefoil: top arc from left-of-tangle going up-and-over to
  % right-of-tangle. This is the "outside" closure of the 2-braid.
  \draw[#3, line width=1.1pt, line cap=round]
    ({\xL-\hw}, \ytop) -- ({\xL-\hw}, {\ytop+0.35})
      arc[start angle=180, end angle=0, x radius=\hw, y radius=0.32]
      -- ({\xL+\hw}, \ytop);
  % Closure of right trefoil: same on the right side
  \draw[#3, line width=1.1pt, line cap=round]
    ({\xR-\hw}, \ytop) -- ({\xR-\hw}, {\ytop+0.35})
      arc[start angle=180, end angle=0, x radius=\hw, y radius=0.32]
      -- ({\xR+\hw}, \ytop);
  % Connect-sum band at the bottom: trivial band joining the two trefoils.
  % It loops from the bottom-right of the left tangle to the bottom-left
  % of the right tangle, dipping down. Then we close off the two outer
  % bottom strands with a large arc going around.
  \draw[#3, line width=1.1pt, line cap=round]
    ({\xL+\hw}, \ybot)
      arc[start angle=180, end angle=360, x radius={(\xR-\xL-\tanglew)/2}, y radius=0.55];
  % Outer bottom closure: leftmost-bottom to rightmost-bottom around the underside.
  \draw[#3, line width=1.1pt, line cap=round]
    ({\xL-\hw}, \ybot) -- ({\xL-\hw}, {\ybot-0.45})
      arc[start angle=180, end angle=360, x radius={(\xR-\xL)/2+\hw}, y radius=0.85]
      -- ({\xR+\hw}, \ybot);
}
 in your dissertation preamble.
%
% Requires the following packages and libraries to be loaded:
%   \usepackage{tikz}
%   \usepackage{xcolor}
%   \usetikzlibrary{calc}
%
% Defines:
%   \crossing{x}{y}{w}{h}{sign}{color}    -- single crossing
%   \tangle{xc}{ytop}{ybot}{k}{color}     -- vertical tangle of k crossings
%   \pretzel{p1}{p2}{p3}{color}           -- 3-tangle pretzel P(p1,p2,p3)
%   \bandhighlight                        -- dashed pink box on third tangle
%
% These also define the colours \thesisblue, \thesisred, \bandfill
% (used in the dissertation's banner/highlight scheme).
% =====================================================================

% --- Colours ---------------------------------------------------------
\providecolor{thesisblue}{HTML}{1F4E79}
\providecolor{thesisred}{HTML}{C00000}
\providecolor{bandfill}{HTML}{FFE5E5}

% ---------------------------------------------------------------------
% \crossing{x}{y}{w}{h}{sign}{color}
%   A single crossing centred at (x, y), of width w and height h.
%   sign > 0: right-handed (SW->NE strand over NW->SE strand).
%   sign < 0: left-handed  (NW->SE strand over SW->NE strand).
% ---------------------------------------------------------------------
\newcommand{\crossing}[6]{%
  \pgfmathsetmacro{\eps}{0.08*#3}%
  \pgfmathsetmacro{\hw}{0.5*#3}%
  \pgfmathsetmacro{\hh}{0.5*#4}%
  \ifnum#5>0\relax
    \draw[#6, line width=1.1pt, line cap=round]
      ({#1-\hw},{#2+\hh}) -- ({#1-\eps},{#2+\eps*#4/#3});
    \draw[#6, line width=1.1pt, line cap=round]
      ({#1+\eps},{#2-\eps*#4/#3}) -- ({#1+\hw},{#2-\hh});
    \draw[#6, line width=1.1pt, line cap=round]
      ({#1-\hw},{#2-\hh}) -- ({#1+\hw},{#2+\hh});
  \else
    \draw[#6, line width=1.1pt, line cap=round]
      ({#1-\hw},{#2-\hh}) -- ({#1-\eps},{#2-\eps*#4/#3});
    \draw[#6, line width=1.1pt, line cap=round]
      ({#1+\eps},{#2+\eps*#4/#3}) -- ({#1+\hw},{#2+\hh});
    \draw[#6, line width=1.1pt, line cap=round]
      ({#1-\hw},{#2+\hh}) -- ({#1+\hw},{#2-\hh});
  \fi
}

% ---------------------------------------------------------------------
% \tangle{xc}{ytop}{ybot}{k}{color}
%   A vertical tangle of |k| crossings of sign(k) between y = ytop
%   and y = ybot, centred horizontally at x = xc. k = 0 gives two
%   parallel vertical strands.
% ---------------------------------------------------------------------
\newcommand{\tangle}[5]{%
  \pgfmathsetmacro{\tanglew}{0.55}%
  \pgfmathsetmacro{\hw}{0.5*\tanglew}%
  \ifnum#4=0\relax
    \draw[#5, line width=1.1pt, line cap=round]
      ({#1-\hw},#2) -- ({#1-\hw},#3);
    \draw[#5, line width=1.1pt, line cap=round]
      ({#1+\hw},#2) -- ({#1+\hw},#3);
  \else
    \ifnum#4>0\relax
      \def\sgnval{1}%
      \pgfmathtruncatemacro{\absk}{#4}%
    \else
      \def\sgnval{-1}%
      \pgfmathtruncatemacro{\absk}{-(#4)}%
    \fi
    \pgfmathsetmacro{\totalh}{#2-(#3)}%
    \pgfmathsetmacro{\hper}{\totalh/\absk}%
    \foreach \i in {1,...,\absk} {%
      \pgfmathsetmacro{\cy}{#2-(\i-0.5)*\hper}%
      \crossing{#1}{\cy}{\tanglew}{\hper*0.85}{\sgnval}{#5}%
    }%
  \fi
}

% ---------------------------------------------------------------------
% \pretzel{p1}{p2}{p3}{color}
%   The standard diagram of the 3-tangle pretzel P(p1, p2, p3),
%   centred at (0, 0) and approximately 7 x 8 units in size.
% ---------------------------------------------------------------------
\newcommand{\pretzel}[4]{%
  \pgfmathsetmacro{\tanglew}{0.55}%
  \pgfmathsetmacro{\hw}{0.5*\tanglew}%
  \def\ytop{1.6}%
  \def\ybot{-1.6}%
  \def\xleft{-2.0}%
  \def\xmid{0.0}%
  \def\xright{2.0}%
  \tangle{\xleft}{\ytop}{\ybot}{#1}{#4}%
  \tangle{\xmid}{\ytop}{\ybot}{#2}{#4}%
  \tangle{\xright}{\ytop}{\ybot}{#3}{#4}%
  \draw[#4, line width=1.1pt, line cap=round]
    ({\xleft+\hw}, \ytop)
      arc[start angle=180, end angle=0, x radius=0.725, y radius=0.32];
  \draw[#4, line width=1.1pt, line cap=round]
    ({\xmid+\hw}, \ytop)
      arc[start angle=180, end angle=0, x radius=0.725, y radius=0.32];
  \draw[#4, line width=1.1pt, line cap=round]
    ({\xleft+\hw}, \ybot)
      arc[start angle=-180, end angle=0, x radius=0.725, y radius=0.32];
  \draw[#4, line width=1.1pt, line cap=round]
    ({\xmid+\hw}, \ybot)
      arc[start angle=-180, end angle=0, x radius=0.725, y radius=0.32];
  \draw[#4, line width=1.1pt, line cap=round]
    ({\xleft-\hw}, \ytop) -- ({\xleft-\hw}, {\ytop+0.4})
      arc[start angle=180, end angle=0,
          x radius={(\xright-\xleft)/2+\hw}, y radius=0.85]
      -- ({\xright+\hw}, \ytop);
  \draw[#4, line width=1.1pt, line cap=round]
    ({\xleft-\hw}, \ybot) -- ({\xleft-\hw}, {\ybot-0.4})
      arc[start angle=-180, end angle=0,
          x radius={(\xright-\xleft)/2+\hw}, y radius=0.85]
      -- ({\xright+\hw}, \ybot);
  \node[anchor=north, font=\bfseries, fill=white, inner sep=2pt]
    at (\xleft,  {\ybot-1.4}) {$#1$};
  \node[anchor=north, font=\bfseries, fill=white, inner sep=2pt]
    at (\xmid,   {\ybot-1.4}) {$#2$};
  \node[anchor=north, font=\bfseries, fill=white, inner sep=2pt]
    at (\xright, {\ybot-1.4}) {$#3$};
}

% ---------------------------------------------------------------------
% \bandhighlight
%   A dashed pink box around the third tangle of \pretzel, highlighting
%   it as the parametric "twisting band" of the Hedden-Watson construction.
% ---------------------------------------------------------------------
\newcommand{\bandhighlight}{%
  \draw[thesisred, dashed, line width=0.8pt,
        fill=bandfill, fill opacity=0.35]
    (1.55, -1.75) rectangle (2.45, 1.75);
}

% ---------------------------------------------------------------------
% \trefoilcs{s1}{s2}{color}
%   The standard diagram of the connect-sum T(s1) # T(s2) of two
%   trefoils, where s1, s2 in {+1, -1} indicate the chirality of each
%   summand (+1 = right-handed, -1 = left-handed).
%
%   Drawn as two side-by-side closed 2-braids joined by the trivial
%   connect-sum band. Centred at (0, 0).
% ---------------------------------------------------------------------
\newcommand{\trefoilcs}[3]{%
  \pgfmathsetmacro{\tanglew}{0.55}%
  \pgfmathsetmacro{\hw}{0.5*\tanglew}%
  \def\ytop{1.5}%
  \def\ybot{-1.5}%
  \def\xL{-1.5}%
  \def\xR{1.5}%
  % Pre-evaluate the crossing counts (3 with the sign of s1, s2)
  \pgfmathtruncatemacro{\trefL}{(#1)*3}%
  \pgfmathtruncatemacro{\trefR}{(#2)*3}%
  % Two trefoil tangles (each = 3 crossings of given chirality)
  \tangle{\xL}{\ytop}{\ybot}{\trefL}{#3}%
  \tangle{\xR}{\ytop}{\ybot}{\trefR}{#3}%
  % Closure of left trefoil: top arc from left-of-tangle going up-and-over to
  % right-of-tangle. This is the "outside" closure of the 2-braid.
  \draw[#3, line width=1.1pt, line cap=round]
    ({\xL-\hw}, \ytop) -- ({\xL-\hw}, {\ytop+0.35})
      arc[start angle=180, end angle=0, x radius=\hw, y radius=0.32]
      -- ({\xL+\hw}, \ytop);
  % Closure of right trefoil: same on the right side
  \draw[#3, line width=1.1pt, line cap=round]
    ({\xR-\hw}, \ytop) -- ({\xR-\hw}, {\ytop+0.35})
      arc[start angle=180, end angle=0, x radius=\hw, y radius=0.32]
      -- ({\xR+\hw}, \ytop);
  % Connect-sum band at the bottom: trivial band joining the two trefoils.
  % It loops from the bottom-right of the left tangle to the bottom-left
  % of the right tangle, dipping down. Then we close off the two outer
  % bottom strands with a large arc going around.
  \draw[#3, line width=1.1pt, line cap=round]
    ({\xL+\hw}, \ybot)
      arc[start angle=180, end angle=360, x radius={(\xR-\xL-\tanglew)/2}, y radius=0.55];
  % Outer bottom closure: leftmost-bottom to rightmost-bottom around the underside.
  \draw[#3, line width=1.1pt, line cap=round]
    ({\xL-\hw}, \ybot) -- ({\xL-\hw}, {\ybot-0.45})
      arc[start angle=180, end angle=360, x radius={(\xR-\xL)/2+\hw}, y radius=0.85]
      -- ({\xR+\hw}, \ybot);
}
